\chapter{Creación de Personajes}

En el sistema \textbf{ZINA}, un personaje no es un compendio de habilidades matemáticas abstractas ni una lista de poderes superheroicos; es un habitante de la Margen Izquierda o de los Zócalos que intenta sobrevivir, proteger a los suyos y conservar un resto de humanidad bajo el peso aplastante del asfalto, el humo y el control corporativo de Orekak-Corp.

Crear un personaje en \textit{NeoBilbao} es un proceso narrativo y ágil que se realiza en grupo antes de encender la forja de la primera sesión. Sigue los siguientes seis pasos:

\begin{figure}[htbp]
  \centering
  \includegraphics[width=\textwidth]{08_arquetipos_la_margen.jpg}
  \caption{Habitantes y cuadrillas de la Margen: el pulso de la resistencia cotidiana.}
\end{figure}

\section{Paso 1: Elegir la Era de Juego}
Antes de repartir las fichas, el grupo debe acordar en qué periodo histórico transcurrirá la crónica:
\begin{itemize}
    \item \textbf{2084 -- La Fractura Temprana:} La ciudad se desangra entre la vieja identidad fabril y el nacimiento del monopolio de Orekak-Corp. La tecnología es pesada, de tubos catódicos, cables coaxiales, humo de carbón y lluvia ácida. Los personajes son supervivientes del colapso, detectives callejeros, hackers de terminal y trabajadores que ven cómo se alza la Cúpula y temen a la Lotería Azul.
    \item \textbf{2102 -- El Siglo Blindado:} Dos décadas después, la Cúpula Dorada brilla sobre los rascacielos oligarcas mientras la clase obrera vive hacinada en los Zócalos de hormigón. El \textit{Velo} biomonitorea cada latido, la aleación de \textit{Burdin-Urdin} fluye por las venas de los trabajadores y el \textit{Etxearen Zina} exige entregar a los hijos como prenda de lealtad colonial.
\end{itemize}

\section{Paso 2: Elegir un Arquetipo}
Cada jugador elige uno de los diez arquetipos disponibles. Cada arquetipo representa un rol social, laboral y táctico dentro de la cuadrilla:

\begin{table}[htbp]
\centering
\small
\begin{tabularx}{\textwidth}{l l X}
\toprule
\textbf{Era} & \textbf{Arquetipo} & \textbf{Esencia y Rol en la Cuadrilla} \\
\midrule
\multirow{5}{*}{\textbf{2084}} 
& \textbf{El Sabueso} & Rastreador urbano, expolicía o cobrador cínico pero leal. \\
& \textbf{La Biohacker} & Cirujana clandestina, artesana de ciberware reciclado. \\
& \textbf{El Celador} & Obrero de choque, custodio vecinal, muro humano. \\
& \textbf{La Burócrata} & Inspectora de archivos, conoce las grietas de la ley corporativa. \\
& \textbf{El Familiar} & Núcleo emocional, buscador incansable de los desaparecidos. \\
\midrule
\multirow{5}{*}{\textbf{2102}} 
& \textbf{La Fundidora} & Piloto de exoesqueleto Basajaun, fuerza bruta industrial. \\
& \textbf{El Técnico} & Mantenedor de conductos, saboteador del Velo y la red. \\
& \textbf{La Delegada} & Líder sindical clandestina, voz de la asamblea y la huelga. \\
& \textbf{El Rebotado} & Paria del Yermo, contrabandista y superviviente salvaje. \\
& \textbf{La Médica} & Facultativa de dispensario fabril, alquimista del plasma. \\
\bottomrule
\end{tabularx}
\caption{Los Diez Arquetipos del ZINA Engine.}
\end{table}

\section{Paso 3: Asignar las Emociones}
En el sistema ZINA no existen atributos físicos convencionales como Fuerza o Agilidad. Lo que define a un habitante de NeoBilbao es su \textbf{disposición emocional}: el impulso anímico bajo el cual reacciona ante la opresión, el peligro y la maquinaria corporativa.

Las seis emociones del sistema son:
\begin{itemize}
    \item \lead{Enfadado:} Furia de clase, frustración acumulada, indignación moral y violencia visceral. Se utiliza cuando el personaje ataca de frente, rompe una puerta a patadas o confronta a un capataz (\textit{Atacar / Forcejear}).
    \item \lead{Asustado:} Instinto de supervivencia, adrenalina pura, terror y prudencia. Se emplea cuando el personaje huye a la desesperada, se oculta en las sombras o esquiva un peligro letal (\textit{Infiltrarse / Ocultarse}).
    \item \lead{Culpable:} El peso de los errores pasados, la vergüenza y el auto-reproche. Es la emoción que impulsa a violar códigos de seguridad, hackear terminales corporativas o asumir riesgos suicidas para enmendar una falta (\textit{Hackear / Puentear Sistemas}).
    \item \lead{Afligido:} Empatía, dolor por la pérdida de seres queridos, tristeza compartida y compasión. Se utiliza para consolar a un herido, tratar traumas médicos o apelar a los sentimientos de otra persona (\textit{Manipular / Conectar}).
    \item \lead{Pacífico:} Calma forzada, desapego analítico, sangre fría y meticulosidad. Se utiliza cuando el personaje opera maquinaria con pulso firme, falsifica un albarán burocrático o analiza un plano técnico sin alterarse (\textit{Analizar con Sangre Fría}).
    \item \lead{Alegre:} Esperanza activa, humor rebelde, calidez humana y orgullo de barrio (\textit{Infundir Esperanza}).
\end{itemize}

\subsection{El reparto de puntos según la Era}
El reparto de modificadores refleja el impacto del paso del tiempo y la opresión en el espíritu obrero:
\begin{itemize}
    \item \textbf{En 2084 (La Margen Izquierda):} La cuadrilla aún conserva el calor humano y la esperanza viva. Reparte los siguientes seis valores entre las seis emociones:
    \[
    \mathbf{+2}, \quad \mathbf{+1}, \quad \mathbf{+1}, \quad \mathbf{0}, \quad \mathbf{0}, \quad \mathbf{-1}
    \]
    \item \textbf{En 2102 (Los Zócalos Industriales):} Tras el Cataclismo, la Guerra del Plasma y el encierro bajo el Velo, la emoción \textbf{Alegre} se ha extinguido (está marcada con una \textbf{X} en la ficha). Reparte los siguientes cinco valores entre las cinco emociones restantes (Enfadado, Asustado, Culpable, Afligido, Pacífico):
    \[
    \mathbf{+2}, \quad \mathbf{+1}, \quad \mathbf{+1}, \quad \mathbf{0}, \quad \mathbf{-1}
    \]
\end{itemize}

\section{Paso 4: Calcular Derivadas y Reservas}
Anota en tu ficha los valores de las reservas vitales del sistema:
\begin{itemize}
    \item \textbf{Zama (Carga Psicológica y Neurológica):} Comienza en \textbf{0} en una escala de \textbf{0 a 6 casillas}. El estrés aumenta al respirar gases ácidos, sufrir sobrecargas de terminales o presenciar la crueldad corporativa. Al marcar la 6.ª casilla, sufres el \textbf{Colapso} (ataque de pánico ciego, síncope o parálisis catatónica).
    \item \textbf{Integridad Física (Heridas):} Se compone de tres umbrales de daño:
    \begin{enumerate}
        \item \textit{Herida Leve:} Magulladuras, cortes superficiales de chapa.
        \item \textit{Herida Grave:} Quemaduras de plasma, fracturas óseas, descarga neural. Aplica $-1$ continuo hasta recibir primeros auxilios.
        \item \textit{Herida Crítica:} Hemorragia severa o fallo multiorgánico. Requiere estabilización inmediata o conduce a la muerte.
    \end{enumerate}
    \item \textbf{Reserva de Auzolan (Solidaridad Obrera):} Al inicio de la sesión, la cuadrilla recibe \textbf{1 punto de Auzolan por cada jugador en la mesa}. Se puede gastar para sumar $+1$ a la tirada de un compañero, absorber daño o declarar un recurso barrial olvidado.
\end{itemize}

\section{Paso 5: Vínculos, Deudas y Juramentos}
Un personaje sin ataduras no tiene sentido en NeoBilbao. Durante la creación, responde en voz alta a las siguientes preguntas con el resto de la mesa:
\begin{enumerate}
    \item \textbf{Vínculo con la Cuadrilla:} Escoge a otro personaje jugador. ¿Por qué le debes la vida o en qué turno de trabajo peligroso aprendisteis a confiar ciegamente el uno en el otro?
    \item \textbf{Deuda o Fisura Moral:} ¿Qué secreto compartes con otro compañero que podría acarrear la deportación o el despido inmediato si los Secretas lo descubrieran?
    \item \textbf{El Juramento Personal (Zina):} Define una promesa inquebrantable. Ejemplo: \textit{«Sacaré a mi sobrina de los barracones de filtrado»}, \textit{«Descubriré quién firmó la orden de desalojo de Sestao»} o \textit{«Haremos que las sirenas del Alto Horno vuelvan a sonar libres»}.
\end{enumerate}

\section{Paso 6: Equipo y Wetware Inicial}
Cada arquetipo otorga un conjunto de herramientas de trabajo, armas improvisadas o reglamentarias y piezas de cibernética inicial. Además, todo personaje comienza con:
\begin{itemize}
    \item Una tarjeta de identificación laboral o ficha de racionamiento con chip magnético.
    \item Una prenda de abrigo impermeable tratada contra la lluvia cáustica (chubasquero de lona o mono de faena ignífugo).
    \item Un objeto personal cargado de memoria: una foto analógica plastificada, un mechero de mecha grabado, una medalla de la Virgen de Begoña oxidada o una cinta magnética con canciones previas a la guerra.
\end{itemize}

\begin{CajaEjemplo}{Ejemplo Completo de Creación: Maitane, Fundidora de 2102}
\textbf{Jugadora:} Leire decide crear a una trabajadora de primera línea en los hornos de Burdin-Urdin.
\begin{itemize}
    \item \textbf{Arquetipo:} La Fundidora (2102).
    \item \textbf{Concepto:} Operaria de colada continua en el Alto Horno 4, expulsada del turno principal por organizar un paro de brazos caídos.
    \item \textbf{Disposición Emocional (2102):} Enfadado $+2$ (furia obrera ante la injusticia), Asustado $+1$ (instinto para no quedar atrapada en el foso térmico), Pacífico $+1$ (pulso firme manejando toneladas de metal), Culpable $0$ (se reprocha no haber evitado la detención de su hermano), Afligido $-1$ (se ha vuelto dura e impermeable al dolor ajeno). \textit{Alegre: Extinguido [X]}.
    \item \textbf{Movimientos elegidos:} \textit{Piel de Amianto} (reduce el daño térmico y de plasma) y \textit{Llave de Trinquete} (bonificación de fuerza al combatir con herramientas pesadas).
    \item \textbf{Equipo:} Exoesqueleto de carga Basajaun-MK3 desgastado, soplete térmico industrial (\textit{Daño 2, Peligrosa, Penetrante}), mono ignífugo con parches de la asamblea y un colgante con la llave del taquillero de su hermano ausente.
    \item \textbf{Juramento:} \textit{«Nadie en mi cuadrilla se quedará atrapado detrás de una compuerta cerrada mientras yo tenga combustible en el soplete».}
\end{itemize}
\end{CajaEjemplo}
